\subsection{Deployment Planner}

The Deployment Planner transforms the results of the preceding stages into a
deployment plan on a shared, heterogeneous infrastructure. Unlike the
analysis stages, which rely on an LLM to interpret the application, the
planner is a deterministic optimization procedure: it consumes the structured
profiles produced by these stages and returns the same plan for the same
inputs, which makes every decision verifiable and allows a direct comparison
with baseline strategies.

\paragraph{Infrastructure model.}
The infrastructure is a graph $\mathcal{I}=(\mathcal{N}\cup\mathcal{S}w,\mathcal{L})$,
where $\mathcal{N}=\{n_1,\ldots,n_k\}$ is the set of compute nodes,
$\mathcal{S}w$ a set of network switches, and $\mathcal{L}$ the set of links.
Each node $n$ belongs to a tier (device, Edge, or Cloud) and to a site with a
region, and is characterized by its available capacity
$\kappa_n=(\kappa_n^{cpu},\kappa_n^{mem},\kappa_n^{gpu},\kappa_n^{sto})$, its
availability $a_n$, and whether its resources are comparable with those of a
general-purpose server (microcontrollers are not). Each link $\ell$ has a
bandwidth $B_\ell$, a latency $\lambda_\ell$, and a type (LAN, Wi-Fi,
Bluetooth, or WAN). Local devices are bound to the anchors identified by the
Offloadability Analysis, and each unit is associated with its \emph{original
host}, i.e., the device on which it executes before planning (which may differ
from the node on which it is eventually deployed).

\paragraph{Decision variables.}
For each execution unit $u_i$ and node $n_j$,
\begin{equation}
x_{ij} =
\begin{cases}
1, & \text{if an instance of } u_i \text{ is deployed on } n_j,\\
0, & \text{otherwise},
\end{cases}
\end{equation}
and $k_i=\sum_j x_{ij}$ is the number of replicas of $u_i$. At most one replica
of a unit is placed on a node (anti-affinity).

\paragraph{Constraints.}
A configuration is feasible when:
\begin{enumerate}
\item \emph{Anchoring}: every non-offloadable unit is placed on the node bound
to its anchor, and units connected by a hardware edge share a node.
\item \emph{Co-location}: every initialization unit that follows other units
is deployed on each node hosting one of the units it supports; only its
persistent storage is reserved, and its start-up peak must fit.
\item \emph{Capacity}: for each node and resource $\rho$,
\begin{equation}
\sum_i x_{ij}\, \mu_{\pi_i}\, r_i^{\rho} + o_j^{\rho} \le \theta\, \kappa_j^{\rho},
\end{equation}
where $r_i^{\rho}$ is the estimated demand, $\mu_{\pi_i}$ a safety margin that
depends on the provenance of the estimate (e.g., $1.0$ for observed and
$1.25$ for assumed values), $o_j^{\rho}$ the runtime overhead of hosting the
application on $n_j$, and $\theta$ the utilization target of the SLA policy.
\item \emph{GPU and services}: units requiring a GPU are placed only on nodes
with sufficient GPU memory, and units requiring an external service only on
nodes with Internet access.
\item \emph{Proximity}: each proximity condition (same site or same region) of
the offloadability profile holds between the nodes hosting the two units.
\item \emph{Bandwidth}: the traffic of each edge $(u_s,u_t)$ is split evenly
among the replica pairs and routed along the lowest-latency path; for each
link, $\sum \mathit{traffic}_\ell \le \theta\, B_\ell$.
\end{enumerate}

\paragraph{Objective.}
Among feasible configurations, the planner minimizes
\begin{equation}
J(c)=\sum_{q} w_q \,\frac{\phi_q(c)}{\nu_q},
\end{equation}
where the components $\phi_q$ are the total inter-node traffic, the traffic
crossing WAN links, a latency cost $\sum_{(s,t)} \omega_{st}\,\lambda_{st}$
weighted by the latency sensitivity $\omega_{st}$ of each flow, the maximum
load of local devices, two load-balancing terms defined below, the
Cloud resources used, the number of offloaded units whose decision depends on
an unverified assumption, and the availability shortfall with respect to the
SLA target, with availability $1-\prod_{j:x_{ij}=1}(1-a_j)$. The weights $w_q$
and normalization constants $\nu_q$ are part of the SLA policy
$P_{\mathrm{SLA}}$.

\paragraph{Load balancing and workload distribution.}
Load balancing is evaluated over the whole pool $\mathcal{P}$ of shared Edge
and Cloud nodes, including idle nodes and the load already reserved by
previously planned applications. For each node, the dominant utilization
\begin{equation}
\eta_j = \max_{\rho}\ \frac{\text{used}_j^{\rho} + \text{load}_j^{\rho}}{\kappa_j^{\rho}}
\end{equation}
is taken over CPU, memory, GPU memory and storage, so that the scarcest
resource of each node (typically GPU memory) determines its load. Two terms
enter the objective: the dispersion $\sigma(\{\eta_j\}_{j\in\mathcal{P}})$,
which penalizes uneven load, and the bottleneck $\max_{j\in\mathcal{P}}\eta_j$,
which penalizes hotspots. Both are measured as the marginal change caused by
the application being planned, since the load of other applications is a
constant that the application cannot influence. Consolidating all units on a
single node therefore no longer yields a spuriously balanced configuration.
Workload is distributed among the replicas of a unit: the traffic of each edge
is split evenly among replica pairs, and the CPU demand of a replicated unit is
divided among its $k$ replicas, each also incurring a fixed per-replica
overhead, whereas memory, GPU memory and storage are reserved in full by every
replica, since each instance loads its own model, state and data. The weights
of the two balancing terms control the trade-off between spreading load and
keeping communicating units close; setting them to zero yields a purely
communication- and latency-driven placement.

\paragraph{Availability semantics.}
The availability of a unit is $A_i = 1-\prod_{j:x_{ij}=1}(1-a_j)$ and the SLA
policy defines a target $A^\ast$. Two cases are distinguished. For anchored
units, $A_i$ is bounded by the availability of the device they depend on and
cannot be changed by any placement decision; these units are therefore
excluded from the availability requirement and are reported separately. For
movable units, the policy selects one of two modes. In \emph{soft} mode (the
default), availability enters the objective through the shortfall term
$\sum_i \max(0, A^\ast - A_i)$, so a plan may be feasible while some units
remain below the target; each such unit is reported together with its cause.
In \emph{hard} mode, the condition $A_i \ge A^\ast$ is added to the
feasibility constraints and checked after replication; a plan that violates it
is rejected. This distinction makes explicit whether an SLA shortfall results
from the infrastructure (anchored devices), from the SLA itself (too few
requested replicas), or from placement constraints that bind a unit to a less
available host.

\paragraph{Procedure.}
The planner first fixes the anchored units, then generates for every movable
unit the candidate nodes that satisfy its unit-level conditions (capacity
type, GPU, external access, proximity to anchored peers). When the candidate
space is small, all configurations are evaluated exhaustively; otherwise a
greedy assignment refined by local search is used. Replication is part of the
optimization rather than a post-processing step: every feasible base
configuration is extended with replicas, which are placed on distinct feasible
nodes (anti-affinity), and the plan with the lowest final objective is
selected. Two replication policies are supported. Under \emph{SLA-required}
replication, the number of replicas requested by $P_{\mathrm{SLA}}$ is always
deployed, even if the additional traffic and load increase $J$; under
\emph{score-driven} replication, a replica is added only if it decreases $J$.
How replicas share the work of a unit is made explicit by its
\emph{replication semantics}, derived from the upstream profiles:
\begin{itemize}
\item \emph{load-split} (stateless units): the incoming traffic and the CPU
work are divided among the $k$ replicas;
\item \emph{mirrored} (units holding persistent data, e.g., an event log or a
historical store): writes are mirrored, i.e., every replica receives all
incoming writes and keeps a full copy, so the write traffic is multiplied by
$k$ and each replica reserves the full CPU and storage demand; reads (outgoing
flows) are load-balanced, each being served by a single replica (optionally,
by a designated primary replica); replicas are assumed to be kept consistent by
asynchronous write replication, whose ordering and consistency guarantees are
not modelled;
\item \emph{partitioned} (optional alternative for persistent units): the data
are sharded, each replica receiving and storing $1/k$ of the records, without
redundancy of individual records;
\item \emph{single}: units that keep in-memory processing state (e.g., object
tracking or sliding frame buffers) are not replicated unless the policy
explicitly allows it.
\end{itemize}
Each semantics is defined in the SLA policy by four rules, applied
consistently to capacity reservation, traffic routing and reporting: the
distribution of incoming writes (all replicas or split), of outgoing reads
(load-balanced or primary), of CPU work (split or full), and of storage (split
or full). Every replication decision is reported together with its cause, its
semantics, and the resulting per-replica shares of write traffic, read
traffic, CPU and storage. Applications are planned in order of priority on the residual
capacity of the shared infrastructure, so that contention between applications
is taken into account. To make every decision traceable, the planner attaches a rationale to each
unit: anchored units report the blocking factor that binds them; follower units
report the units they are co-located with; and each movable unit reports the
nodes excluded before the search (with the reason), and a counterfactual
analysis in which the unit alone is moved to every other candidate node while
the rest of the plan is kept, giving either the resulting change in $J$ and
the objective terms responsible for it, or the constraint the move would
violate. Violations name the units or flows responsible (e.g., the flows that
saturate a link). For evaluation, the resulting plan is compared with
three baselines: keeping all movable units on their original host, placing all
of them in the Cloud, and a greedy first-fit placement.

\begin{algorithm}[ht]
\caption{Deployment Planner}
\label{alg:deployment_planner}
\begin{algorithmic}[1]
\Require Execution units $\mathcal{U}$, resource profile $R$, offloadability profile $O$
\Require Infrastructure $\mathcal{I}$ (residual capacity), SLA policy $P_{\mathrm{SLA}}$
\Ensure Deployment descriptor $D_{\mathrm{deploy}}$
\State $F \leftarrow$ place each non-offloadable unit on the node of its anchor
\State $\mathcal{U}_m \leftarrow$ units that are neither anchored nor followers
\For{each $u_i \in \mathcal{U}_m$}
  \State $N_i \leftarrow \mathrm{GenerateCandidates}(u_i, R, O, \mathcal{I})$
\EndFor
\State $C \leftarrow N_1\times\cdots\times N_{|\mathcal{U}_m|}$ \Comment{or greedy + local search if $|C|$ is too large}
\State $C_{\mathrm{valid}} \leftarrow \emptyset$
\For{each configuration $c \in C$}
  \State $c \leftarrow c \cup F \cup \mathrm{CoLocateFollowers}(c)$
  \If{$\mathrm{SatisfiesResources}(c,R,\mathcal{I})$ \textbf{and} $\mathrm{SatisfiesOffloadability}(c,O,\mathcal{I})$}
    \State $C_{\mathrm{valid}} \leftarrow C_{\mathrm{valid}} \cup \{(c, J(c))\}$
  \EndIf
\EndFor
\State $C_{\mathrm{rep}} \leftarrow \{\mathrm{Replicate}(c, P_{\mathrm{SLA}}) : (c,\cdot)\in C_{\mathrm{valid}}\}$ \Comment{anti-affinity, stateful units kept single}
\State $c^\ast \leftarrow \arg\min_{c\in C_{\mathrm{rep}},\ c\ \mathrm{feasible}} J(c)$
\State Distribute the traffic of each edge evenly among replica pairs
\State Reserve the resources of $c^\ast$ in $\mathcal{I}$
\State $D_{\mathrm{deploy}} \leftarrow (c^\ast, J(c^\ast), \text{constraint checks}, \text{baselines})$
\State \Return $D_{\mathrm{deploy}}$
\end{algorithmic}
\end{algorithm}
